\section{Discussion}
\label{sec:discussion}

% TODO: structure as 3 paragraphs.

% ¶1 --- summarize the three contributions and how they interlock.
%   - SO(2) equivariance: accuracy at moderate cost.
%   - Attention: longer-range correlations within each MP step.
%   - Native ZBL: extends domain of validity to short interatomic
%     distances without breaking conservativeness.

% ¶2 --- limitations.
%   - DPA4 is presented here as a single-task / per-dataset model; we
%     have not yet integrated DPA3's multi-task dataset-encoding
%     scheme, nor trained a DPA4-based foundation potential. We
%     anticipate doing this in follow-up work.
%   - Charged species, open-shell magnetism, and relativistic heavy
%     elements remain out of scope.
%   - SO(2) is in principle a strict subset of SO(3) equivariance; the
%     gap appears empirically negligible but may resurface in
%     symmetry-sensitive tasks (e.g. accurate phonons in highly
%     symmetric crystals).

% ¶3 --- outlook.
%   - DPA4 as the base architecture for a future foundation potential
%     trained on OpenLAM-v1 / OMol25 in a multi-task setup.
%   - Downstream applications enabled by DPA4's efficiency: long-time
%     MD, irradiation-damage simulation, high-throughput catalysis
%     screening on commodity GPUs.

% =============================================================
